% Mathematical sources: mathematics/current/skill/research/20260912.tex,
% lines 175--220; mathematics/current/curvature_box.md, first two sections;
% mathematics/current/rigidity.md, Sections 1--3 (lines 8--235).
\section{Pointwise Ricci pinching and spectral exclusion}
\label{sec:early-pointwise-rigidity}

The integral obstruction has two useful pointwise precursors.  The first
retains the correlation between the four Ricci eigenvalues in the
Chern--Weil density and proves strictness at its algebraic endpoint.
The second uses the harmonic-Weyl equation to exclude a fixed spectral
box.  Its proof also identifies an obstruction to extending the same
scalar-multiplier strategy over the whole curvature space.

\subsection{The strict ratio-three endpoint}

Write $E=\Ric^\circ=\Ric-sg/4$.  Ricci tensors have their tensor norm,
and Weyl tensors have their Hilbert--Schmidt endomorphism norm.  Thus
\begin{equation}\label{eq:epr-cw}
 4\pi^2(2\chi\pm3\tau)
 =\int_M\left(2|W^\pm|^2+\frac{s^2}{24}
                                      -\frac{|E|^2}{2}\right)d\mu.
\end{equation}
For $\sigma\in\{1,-1\}$, set
\[
 \delta_\sigma(g)
 =\|g^{-1}\Ric_g-\sigma I\|_{L^\infty,\mathrm{op}}.
\]

\begin{proposition}[Strictness at Ricci eigenvalue ratio three]
\label{prop:epr-ratio-three}
Let $M$ be a closed connected oriented four-manifold.  Suppose
$\sigma\Ric_g$ is positive definite and, at every point,
\begin{equation}\label{eq:epr-ratio-three}
 \lambda_{\max}(\sigma g^{-1}\Ric_g)
 \leq3\lambda_{\min}(\sigma g^{-1}\Ric_g).
\end{equation}
Then $2\chi(M)>3|\tau(M)|$.  In particular, the conclusion holds if
$\delta_\sigma(g)\leq1/2$.
\end{proposition}

\begin{proof}
Let $r_1,\ldots,r_4>0$ be the eigenvalues of
$\sigma g^{-1}\Ric_g$, and put
\[
 a=\min_i r_i,\qquad \mathsf S=\sum_i r_i,
 \qquad \mathsf Q=\sum_i r_i^2.
\]
The inequalities $a\leq r_i\leq3a$ imply
$r_i^2\leq4ar_i-3a^2$.  Consequently
\begin{equation}\label{eq:epr-density-endpoint}
 \mathsf Q\leq4a\mathsf S-12a^2,\qquad
 F:=\frac{s^2}{24}-\frac{|E|^2}{2}
 =\frac{\mathsf S^2-3\mathsf Q}{6}
 \geq\frac{(\mathsf S-6a)^2}{6}\geq0.
\end{equation}
Equality $F=0$ holds precisely when the eigenvalues are a permutation
of $(3a,a,a,a)$.  Indeed, equality requires each $r_i$ to be $a$ or
$3a$ and their sum to be $6a$.

Equation~\eqref{eq:epr-cw} first gives $2\chi\geq3|\tau|$.
Suppose equality holds.  Then $F\equiv0$ and one Weyl half vanishes.
The simple eigenvalue $3\sigma a$ of $g^{-1}\Ric_g$ defines a smooth
real eigenline bundle.  On its orientation cover, of degree at most
two, choose a smooth unit section $V$.  This nowhere-zero vector
field makes the Euler characteristic of the cover zero; multiplicativity
under finite covers gives $\chi(M)=0$.  The assumed equality then
gives $\tau(M)=0$.  Applying \eqref{eq:epr-cw} to the other chirality
shows that the remaining Weyl half also vanishes.

Work on that closed cover and put $\theta=V^\flat$.  The equality
eigenvalues give
\[
 \Ric=\sigma a(g+2\theta\otimes\theta),\qquad
 s=6\sigma a,
 \qquad
 \mathsf P=\tfrac12(\Ric-sg/6)
             =\sigma a\,\theta\otimes\theta,
\]
where $\mathsf P$ is the Schouten tensor.  Since $W=0$,
$\mathsf P$ is Codazzi.  For $X,Y\perp V$, comparison of
$(\nabla_X\mathsf P)(V,Y)$ and
$(\nabla_V\mathsf P)(X,Y)$ gives
$\langle\nabla_XV,Y\rangle=0$.  Since $V$ has unit length,
$\nabla_XV=0$ for every $X\perp V$.

With $B=\nabla_VV\perp V$, it follows that
\begin{equation}\label{eq:epr-eigenline}
 \nabla V=\theta\otimes B,\qquad
 d\theta=\theta\wedge B^\flat,\qquad
 d^*\theta=0,\qquad
 |d\theta|^2=|\nabla\theta|^2.
\end{equation}
The integrated Bochner identity for this one-form is
\[
 \int\bigl(|d\theta|^2+|d^*\theta|^2\bigr)
 =\int\bigl(|\nabla\theta|^2+\Ric(V,V)\bigr).
\]
Thus $\int\Ric(V,V)=0$, whereas $\Ric(V,V)=3\sigma a$ has a
strict fixed sign.  This contradiction proves strictness.  Finally,
$\delta_\sigma\leq1/2$ places the $r_i$ in $[1/2,3/2]$, which
implies \eqref{eq:epr-ratio-three}.
\end{proof}

Hence every metric on a manifold with $2\chi\leq3|\tau|$ has
$\delta_\sigma>1/2$.  This pointwise strict inequality for each
metric does not assert that the infimum over all metrics is strictly
larger than $1/2$.

\begin{lemma}[The sharp Ricci-box density]
\label{lem:epr-ricci-box}
For $0\leq\delta\leq1/2$, the exact minimum of $F$ over Ricci
eigenvalues satisfying $\delta_\sigma\leq\delta$ is
\begin{equation}\label{eq:epr-box-minimum}
 q_*(\delta)=
 \begin{cases}
  \dfrac23(1-\delta)^2,&0\leq\delta\leq1/3,\\[1mm]
  \dfrac23(1+\delta)(1-2\delta),&1/3\leq\delta\leq1/2.
 \end{cases}
\end{equation}
Consequently
\begin{equation}\label{eq:epr-box-topology}
 4\pi^2(2\chi\pm3\tau)\geq q_*(\delta)\Vol_g(M).
\end{equation}
The minimum is algebraically sharp; no closed metric realizing a
minimizing curvature pattern everywhere is asserted.
\end{lemma}

\begin{proof}
Use $r_i\in[1-\delta,1+\delta]$ and the expression for $F$ in
\eqref{eq:epr-density-endpoint}.  Its second derivative in each
individual coordinate $r_i$, with the others fixed, is $-2/3$.
Repeated minimization in one coordinate therefore places a minimum
at a vertex of the box.  If $m$ coordinates equal $1+\delta$, then
\[
 F_m(\delta)=\frac23+\frac{(2m-4)\delta}{3}
       +\left(\frac{(2m-4)^2}{6}-2\right)\delta^2.
\]
The five values are
\begin{align*}
 F_0&=\tfrac23(1-\delta)^2,
 &F_1&=\tfrac23(1+\delta)(1-2\delta),\\
 F_2&=\tfrac23-2\delta^2,
 &F_3&=\tfrac23(1-\delta)(1+2\delta),\\
 F_4&=\tfrac23(1+\delta)^2.
\end{align*}
Now
\[
 F_1-F_0=\tfrac23\delta(1-3\delta),\qquad
 F_2-F_1=\tfrac23\delta(1-\delta),\qquad
 F_3-F_1=\tfrac43\delta,\qquad F_4\geq F_0.
\]
These comparisons prove \eqref{eq:epr-box-minimum}; integration in
\eqref{eq:epr-cw} proves \eqref{eq:epr-box-topology}.
\end{proof}

The density is positive below $\delta=1/2$, and vanishes at that
endpoint.  Proposition~\ref{prop:epr-ratio-three} supplies the
differential argument needed for strict topology at the endpoint.
The number $1/2$ is the sharp algebraic nonnegativity endpoint for
this Ricci box, not an Einstein-existence threshold.

\subsection{The harmonic-Weyl equation and a spectral box}

In the rest of this section we use
$\Delta=\operatorname{tr}_g\nabla^2$, so that
$\int f\Delta f=-\int|df|^2$.  This sign is opposite to the
nonnegative scalar Laplacian used in some later sections.  The
connection trace Laplacian has the same sign.  Put
\[
 A=W^+,\quad x=|A|^2,\quad y=\det A,\qquad
 A_0=\operatorname{diag}(-1/6,1/12,1/12),
 \quad x_0=1/24,\quad y_0=-1/864.
\]
The associated four-index Weyl tensor has squared norm $4x$.
For $\operatorname{div}W^+=0$, the classical harmonic-Weyl formulas
in these conventions are
\begin{equation}\label{eq:epr-weyl-equation}
 \Delta A=\frac{s}{2}A-6A^2+2|A|^2I,
 \qquad
 \frac12\Delta x=|\nabla A|^2+\frac{s}{2}x-18y.
\end{equation}
We use their eigenframe form as recalled in
\href{https://arxiv.org/pdf/1903.11818}{Peng Wu, Section 2}.

At $s=-4$ and $A=A_0$, the scalar reaction in the second identity is
\begin{equation}\label{eq:epr-norm-reaction}
 \frac{s}{2}x_0-18y_0=-\frac1{12}+\frac1{48}=-\frac1{16}.
\end{equation}
Also
\begin{equation}\label{eq:epr-matrix-reaction}
 A_0^2=-\frac1{12}A_0+\frac1{72}I,
 \qquad -2A_0-6A_0^2+2x_0I=-\frac32A_0.
\end{equation}
Thus the ordinary norm identity at this algebraic spectrum calls for
gradient energy $1/16$.  Even if one could pass its cubic term to a
limit, it would not by itself give a contradiction.  The first-order
Bianchi equations exclude an exact constant-spectrum metric, but
small averaged spectral error alone does not remove possible
derivative energy.

\begin{proposition}[An explicit pointwise spectral exclusion]
\label{prop:epr-spectral-box}
Let $M$ be a closed oriented four-manifold with $s=-4$ and
$\operatorname{div}W^+=0$.  There is a point where at least one
ordered eigenvalue of $W^+$ differs from
$(-1/6,1/12,1/12)$ by more than $10^{-4}$.  In particular, this
holds for every closed Einstein four-manifold with $\Ric=-g$.
No volume, diameter, injectivity-radius or topological hypothesis
is required.
\end{proposition}

\begin{proof}
Work first on the open dense set where eigenvalue multiplicities
are locally constant.  Let $\lambda_1\leq\lambda_2\leq\lambda_3$
be the ordered eigenvalues.  In a local self-dual eigenframe
$\omega_i$ with $|\omega_i|^2=2$, write the connection one-forms
as $a,b,c$ with the harmonic-Weyl convention, and set
\[
 X=(\lambda_1-\lambda_2)\iota_{a^\sharp}\omega_3,\quad
 Y=(\lambda_2-\lambda_3)\iota_{b^\sharp}\omega_1,\quad
 Z=(\lambda_3-\lambda_1)\iota_{c^\sharp}\omega_2.
\]
Here $X,Y,Z$ are one-forms; the musical isomorphism identifies them
with the corresponding gradient vectors.  The eigenframe equations
give
\begin{equation}\label{eq:epr-first-jets}
 d\lambda_1=-X+Z,\qquad
 d\lambda_2=X-Y,\qquad
 d\lambda_3=Y-Z,
\end{equation}
and
\begin{align}
 \Delta x={}&8(|X|^2+|Y|^2+|Z|^2)
 -4(\langle X,Y\rangle+\langle X,Z\rangle+\langle Y,Z\rangle)
 +sx-36y,\label{eq:epr-delta-x}\\
 \Delta y={}&-4\lambda_3|X|^2-4\lambda_1|Y|^2-4\lambda_2|Z|^2
 -4\lambda_2\langle X,Y\rangle-4\lambda_1\langle X,Z\rangle
 \nonumber\\
 &\hspace{18mm}-4\lambda_3\langle Y,Z\rangle
 +\frac32sy-x^2.\label{eq:epr-delta-y}
\end{align}
Choose the smooth invariant polynomial
\begin{equation}\label{eq:epr-polynomial}
 f(x,y)=x+36y+24(x-1/24)^2,
 \qquad a_0(x)=f_x=1+48(x-1/24).
\end{equation}
For the following calculation put
$p=\lambda_1-\lambda_2$, $q=\lambda_2-\lambda_3$,
and $r=\lambda_3-\lambda_1$.  Then
\[
 dx=2(pX+qY+rZ).
\]
At $s=-4$, equations \eqref{eq:epr-delta-x}--\eqref{eq:epr-delta-y}
and the chain rule yield
\begin{equation}\label{eq:epr-polynomial-laplacian}
 \Delta f=\sum_{i,j=1}^3M_{ij}\langle V_i,V_j\rangle
                       +\mathcal P(x,y),\qquad
 (V_1,V_2,V_3)=(X,Y,Z),
\end{equation}
where $M$ is symmetric and
\begin{align*}
 M_{11}&=8a_0-144\lambda_3+192p^2,
 &M_{22}&=8a_0-144\lambda_1+192q^2,\\
 M_{33}&=8a_0-144\lambda_2+192r^2,
 &M_{12}&=-2a_0-72\lambda_2+192pq,\\
 M_{13}&=-2a_0-72\lambda_1+192pr,
 &M_{23}&=-2a_0-72\lambda_3+192qr,
\end{align*}
and
\begin{equation}\label{eq:epr-polynomial-reaction}
 \mathcal P=(-4x-36y)a_0-216y-36x^2.
\end{equation}
At the target,
\begin{equation}\label{eq:epr-target-matrix}
 M_0=\begin{pmatrix}8&-8&-2\\-8&32&-8\\-2&-8&8\end{pmatrix},
 \qquad \mathcal P_0=\frac1{16}.
\end{equation}
The leading principal minors of $M_0-I$ are $7,153,243$, so
$M_0>I$.

Suppose every eigenvalue stays within
$\varepsilon=10^{-4}$ of its target.  Expansion of $x$ and $y$ gives
\begin{align*}
 |x-x_0|&\leq\tfrac23\varepsilon+3\varepsilon^2
                                      \leq\varepsilon,\\
 |y-y_0|&\leq\tfrac5{144}\varepsilon
              +\tfrac13\varepsilon^2+\varepsilon^3
                                      \leq\varepsilon/20.
\end{align*}
Hence $|a_0-1|\leq48\varepsilon$.
Each diagonal entry of $M-M_0$ is bounded in absolute value by
$720\varepsilon+768\varepsilon^2\leq721\varepsilon$;
each off-diagonal entry is bounded by
$360\varepsilon+768\varepsilon^2\leq361\varepsilon$.
The maximum absolute row sum therefore gives
\begin{equation}\label{eq:epr-matrix-margin}
 \|M-M_0\|_{\mathrm{op}}\leq1443\varepsilon<\frac12,
 \qquad M\geq\frac12I.
\end{equation}
Moreover
\[
 |-4x-36y+1/8|\leq\tfrac{29}{5}\varepsilon,
 \qquad |a_0|<2,
\]
and expansion of \eqref{eq:epr-polynomial-reaction} yields
\[
 |\mathcal P-\mathcal P_0|
 \leq\tfrac{157}{5}\varepsilon+36\varepsilon^2
 <32\varepsilon.
\]
Thus $\mathcal P>1/32$, and
\begin{equation}\label{eq:epr-positive-laplacian}
 \Delta f\geq\tfrac12(|X|^2+|Y|^2+|Z|^2)+\tfrac1{32}
 \geq\tfrac1{32}.
\end{equation}
Although the local eigenframes need not extend across multiplicity
changes, $f$ is a globally smooth polynomial in $W^+$.
The final inequality extends by continuity from the dense
multiplicity strata.  Its integral contradicts $\int_M\Delta f=0$.
\end{proof}

This proposition excludes uniform closeness to a spectral box.  Small
averaged density need not place every point in that box.  Outside it,
the derivative quadratic form of this polynomial need not be positive;
the polynomial also has fourth-order curvature growth.  These features
matter when passing from pointwise exclusion to an integral theorem.

\subsection{An obstruction to global invariant scalar multipliers}

For a $C^2$ function $f$ of a trace-free symmetric matrix, write the
chain rule for the harmonic-Weyl equation as
\[
 \Delta(f(A))=\mathcal Q_f(\nabla A)
              +\mathcal P_f(A),\qquad
 \mathcal P_f(A)=Df_A\left[\frac{s}{2}A-6A^2+2|A|^2I\right].
\]
The derivative term is restricted to first jets satisfying the
algebraic first-order harmonic-Weyl relations.  Positivity on all such
jets is stronger than positivity along one prescribed metric.

\begin{proposition}[Global scalar-multiplier obstruction]
\label{prop:epr-global-multiplier}
Let $f$ be a $C^2$, orthogonally invariant real function on the entire
vector space of trace-free symmetric $3\times3$ matrices.  If
$\mathcal Q_f$ is nonnegative on every algebraically admissible
harmonic-Weyl first jet, then, for scalar curvature $s=-4$,
\[
 \mathcal P_f(A_0)\leq0.
\]
\end{proposition}

\begin{proof}
Restrict to the negative repeated-eigenvalue ray
\[
 A(u)=\frac{u}{\sqrt6}\operatorname{diag}(-2,1,1),
 \qquad u>0,\qquad \psi(u)=f(A(u)).
\]
At this spectrum use the admissible first jet $Y=0$, $Z=-X$.
Equation~\eqref{eq:epr-first-jets} gives
\[
 d\lambda_1=-2X,\qquad d\lambda_2=X,\qquad d\lambda_3=X,
 \qquad |du|^2=6|X|^2,
 \qquad |\nabla A|^2=10|X|^2.
\]
The last equality also follows by comparing
\eqref{eq:epr-delta-x} with \eqref{eq:epr-weyl-equation}.
The derivative of $A$ has a radial part and a part tangent to its
orthogonal-conjugation orbit.  Isotropy at the repeated eigenvalue
makes $Df$ radial.  On the ray and its conjugation orbits,
$f(A)-\psi(|A|)$ and its first derivative vanish; its Hessian
therefore vanishes on the chosen tangent jet.  The derivative term
equals that of the radial function $\psi(|A|)$:
\begin{align}
 \mathcal Q_f
 &=\psi''(u)|du|^2
       +\frac{\psi'(u)}{u}\bigl(|\nabla A|^2-|du|^2\bigr)
       \nonumber\\
 &=\left(\psi''(u)+\frac{2}{3u}\psi'(u)\right)|du|^2.
 \label{eq:epr-radial-derivative}
\end{align}
Nonnegativity for this jet at every $u>0$ implies
\begin{equation}\label{eq:epr-monotone-multiplier}
 (u^{2/3}\psi'(u))'\geq0.
\end{equation}
Smooth invariance at the zero matrix forces $Df_0=0$, since the
only invariant linear functional on trace-free symmetric matrices
is zero.  Thus $\psi'(u)=O(u)$ as $u\downarrow0$.
Integrating \eqref{eq:epr-monotone-multiplier} from zero gives
$\psi'(u)\geq0$ for all $u>0$.

At $u_0=1/\sqrt{24}$, equation~\eqref{eq:epr-matrix-reaction}
gives
\[
 \mathcal P_f(A_0)=-\frac32u_0\psi'(u_0)\leq0.
\]
This proves the assertion using only admissible algebraic first
jets; it does not assume that every such jet extends to a closed
Einstein metric.
\end{proof}

The obstruction allows arbitrary nonhomogeneous growth of $f$.
For the successful local polynomial \eqref{eq:epr-polynomial},
the coefficient in \eqref{eq:epr-radial-derivative} tends to
$-10/3$ as $u\downarrow0$.  A cutoff does not remove the derivative
cost created by this loss of sign.

\begin{corollary}[The local homogeneous obstruction]
\label{cor:epr-homogeneous-multiplier}
No $C^2$, orthogonally invariant, degree-two homogeneous function
near the nonzero target orbit can have both a nonnegative full
derivative quadratic form on all admissible harmonic-Weyl first
jets and a positive algebraic reaction at $A_0$ when $s=-4$.
\end{corollary}

\begin{proof}
Write the invariant homogeneous function locally as
\[
 f(A)=x\varphi(z),\qquad z=\frac{y}{x^{3/2}},
 \qquad z_0=-\frac1{3\sqrt6}.
\]
Euler homogeneity and \eqref{eq:epr-matrix-reaction} give
\begin{equation}\label{eq:epr-homogeneous-reaction}
 \mathcal P_f(A_0)=-3f(A_0)=-\frac18\varphi(z_0).
\end{equation}
Use again $Y=0$, $Z=-X$.  Its radial eigenvalue derivatives are
$(-2X,X,X)$, while the eigenframe connection supplies the
orbit-tangent part of $\nabla A$.  Hence
\[
 |d|A||^2=6|X|^2,\qquad |\nabla A|^2=10|X|^2.
\]
At the repeated eigenvalue, the first matrix derivative of $z$
vanishes, and its Hessian annihilates this radial-and-orbit tangent
jet.  The derivative term consequently is
\begin{equation}\label{eq:epr-homogeneous-derivative}
 \mathcal Q_f=\frac{10}{3}\varphi(z_0)|d|A||^2
                         =20\varphi(z_0)|X|^2.
\end{equation}
Its nonnegativity requires $\varphi(z_0)\geq0$, whereas
\eqref{eq:epr-homogeneous-reaction} is positive only if
$\varphi(z_0)<0$.  These conditions are incompatible.
\end{proof}

The corollary uses homogeneity near the target instead of regularity
at the zero matrix.  It leaves open nonhomogeneous local, singular,
nonlocal and coupled multipliers.  Inverse powers of $|A|$, for
example, alter the radial sign calculation but introduce a zero-set
problem.  None of these pointwise conclusions supplies the
unrestricted averaged weak zero-Weyl endpoint.
